\section{The EventGraph}
\label{sec:eventgraph}

This section introduces the EventGraph: the finite,
visibility-aware directed acyclic graph that all Vegas backends
read.  The graph is the structural core of the language; every
backend-facing notion --- front-running mitigation, information-set
construction, bail-out handling --- is expressed as a property of
this graph or as a rewrite on it.  The frontend pipeline below
treats it as the only stable interface between source and
artifacts.

\begin{center}
\begin{tikzpicture}[
  every node/.style={font=\small},
  box/.style={draw, rounded corners, inner sep=4pt, minimum width=2.2cm,
              minimum height=6mm, fill=blue!6, align=center},
  ir/.style ={draw, rounded corners, inner sep=4pt, minimum width=2.2cm,
              minimum height=6mm, fill=yellow!20, align=center},
  arr/.style={-{Stealth[length=2mm]}, thick}
]
\node[box] (src) at (0,0) {Source\\\TT{.vg}};
\node[box] (ast) at (2.6,0) {AST};
\node[box] (mac) at (5.2,0) {Macro-\\inlined AST};
\node[box] (chk) at (7.8,0) {Type-\\checked};
\node[ir]  (eg)  at (10.4,0) {EventGraph};

\node[box] (sol) at (3.4,-1.6) {Solidity / Vyper};
\node[box] (gam) at (6.0,-1.6) {Gambit EFG};
\node[box] (smt) at (8.4,-1.6) {SMT-LIB};
\node[box] (scr) at (10.7,-1.6) {Scribble};
\node[box] (gal) at (3.4,-2.4) {Gallina / Lean};
\node[box] (mai) at (6.0,-2.4) {MAID JSON};
\node[box] (btc) at (8.4,-2.4) {Bitcoin / LN};
\node[box] (gv)  at (10.7,-2.4) {GraphViz};

\draw[arr] (src) -- (ast);
\draw[arr] (ast) -- (mac);
\draw[arr] (mac) -- (chk);
\draw[arr] (chk) -- (eg);

\foreach \t in {sol, gam, smt, scr, gal, mai, btc, gv}
  \draw[arr, gray] (eg.south) to[bend left=8] (\t.north);
\end{tikzpicture}
\end{center}

The frontend (\file{vegas/frontend/}) parses with an
ANTLR-generated parser, inlines macros, and runs a four-pass type
checker that validates type aliases, macro signatures, the
protocol's join/yield/reveal/commit/withdraw structure, and the
visibility discipline of guards.  It then lowers to an
intermediate representation \TT{GameIR}, which packages the
EventGraph together with a per-role payoff expression and the
chance-role set.  Two transformation passes operate on the IR: the
first builds the EventGraph from typechecked AST (described in
this section), the second is the commit-reveal expansion pass
(\S\ref{sec:commit-reveal}) that rewrites concurrent public writes
into commit-then-reveal pairs.  Backends consume the post-rewrite
graph.

\subsection{Nodes}

A node of the EventGraph represents one statement in the protocol:
a \TT{join}, a \TT{yield}, a \TT{commit}, a \TT{reveal}, or a
\TT{random}.  Each node is identified by a pair
$\mathit{NodeId} = (\Role, n)$ where $\Role$ is the actor and $n$
is an integer giving the node's source position.  The pair is
stable: it is fixed at IR construction time and used throughout
backends as a deterministic identifier.

To each node the IR attaches a payload of three pieces:

\begin{itemize}\itemsep1pt
\item the \emph{specification} (\TT{NodeSpec}): the parameter list,
  any deposit (for \TT{join} nodes), and the guard expression;
\item the \emph{structural metadata} (\TT{NodeStruct}): the owner,
  the set of written fields, the visibility of each written field
  (one of \COMMIT, \REVEAL, \PUBLIC), and the set of fields the
  guard reads from prior nodes;
\item a derived \emph{kind} (\TT{Visibility}), which is the
  ``dominant'' visibility of the node's writes ---
  \COMMIT for hidden writes, \REVEAL for matching reveals,
  \PUBLIC for plain writes.
\end{itemize}

The split between \TT{NodeSpec} and \TT{NodeStruct} is mainly an
engineering choice: backends that need only structural information
(the Gambit and Scribble backends, primarily) can ignore the spec.

\subsection{Edges and dependency inference}
\label{sec:eventgraph:edges}

Edges in the EventGraph represent control or data dependencies.
The IR builds them by analysing the typechecked AST.  An edge
$A \to B$ means $B$ cannot fire until $A$ has fired.  The
compiler installs an edge $A \to B$ when any of the following
holds:

\begin{itemize}\itemsep1pt
\item \emph{Phase order.}  If $A$ and $B$ are in adjacent source
  phases (consecutive top-level statements), $A \to B$.  This
  preserves the textual ordering at the granularity of phases,
  which is what the programmer can rely on.
\item \emph{Guard read.}  If $B$'s guard reads a field $f$, and
  the latest writer of $f$ before $B$ is $A$, then $A \to B$.
  The guard cannot evaluate until the value is available.
\item \emph{Commit before reveal.}  If $B$ is a \REVEAL of field
  $f$ and $A$ is the \COMMIT of $f$, then $A \to B$.  The reveal
  must come after its commit.
\end{itemize}

Within a single phase, multiple roles can act simultaneously (they
appear in the same \TT{yield}/\TT{commit} statement).  Their nodes
share predecessors and are not connected to each other: they are
concurrent in the dependency graph, which is what justifies
treating them as simultaneous moves in the strategic analysis.

\subsection{Visibility}
\label{sec:eventgraph:visibility}

Each written field has a visibility tag that records what an
observer learns from the corresponding write.  The three tags are:

\begin{description}\itemsep1pt
\item[\PUBLIC.]  The value is public from the moment of the write.
  Any later node may read it.
\item[\COMMIT.]  The value is hidden.  Operationally, only its
  cryptographic commitment is public; the actor knows the value;
  no other actor learns it from this node.
\item[\REVEAL.]  A previously committed value is now public.  The
  reveal is paired with the unique prior commit of the same field.
\end{description}

The IR ensures that visibility is consistent with the dependency
relation.  In particular:

\begin{itemize}\itemsep1pt
\item every \REVEAL is reachable from its matching \COMMIT;
\item every field referenced by a guard at node $B$ has either
  been written \PUBLIC by some predecessor of $B$, or has been
  revealed by some predecessor of $B$ (so its value is known
  publicly at $B$);
\item the owner of a node always sees their own writes (an actor
  may guard their action on a field they have just decided ---
  this is how self-only guards like \TT{where x != 2} are
  modelled).
\end{itemize}

If any of these conditions fails, the EventGraph constructor
refuses to build the graph; the compiler surfaces this as a
static error.  Backends therefore work in a setting where the
information flow is fully consistent.

\subsection{Configurations and frontier}
\label{sec:eventgraph:frontier}

A \emph{configuration} of the EventGraph is a partial assignment
from nodes to outcomes (one outcome per node: a tuple of values
for the node's written fields, or \texttt{Quit}).  A node is
\emph{enabled} in a configuration when all its predecessors have
been assigned and its guard, evaluated on the configuration's
public writes plus the actor's local view, is true.  The set of
enabled, unassigned nodes is the configuration's \emph{frontier}.

Two properties of the frontier matter for the backends:

\begin{itemize}\itemsep1pt
\item \emph{Concurrency.}  If two distinct nodes $A, B$ are on the
  same frontier --- neither reaches the other in the DAG ---
  scheduling $A$ before $B$ versus $B$ before $A$ produces the
  same configuration after both have fired.  This is what allows
  Solidity to accept the actions in any order and Gambit to model
  them as simultaneous.
\item \emph{Information barrier.}  If a node $B$ is a \REVEAL,
  every node on which it depends has already fired by the time
  $B$ is enabled.  In particular, in any configuration where
  $B$'s match \COMMIT is unfinished, $B$ is not on the frontier.
  This is the structural counterpart to the commit-reveal
  insertion pass: once a player can reveal, every other player
  has already committed; no player can condition a commit on
  another's revealed value.
\end{itemize}

The compiler computes the frontier on demand with a small
\TT{FrontierMachine} that maintains the set of unresolved nodes
along with the current frontier; resolving the frontier advances
both sets in one step.  The Gambit backend explores the resulting
state space depth-first (\S\ref{sec:gambit}), treating each
frontier resolution as one strategic round.

\subsection{A worked graph}
\label{sec:eventgraph:example}

The Monty Hall example from \S\ref{sec:intro:first-example}
produces the EventGraph below.  Every statement in that source is
its own top-level phase, so the dependency relation is a chain
plus one long-distance edge from the Host's \COMMIT to the
matching \REVEAL.  Squares are \COMMIT writes, hexagons are
\REVEAL writes, ovals are \PUBLIC writes.  The commit-reveal
expansion pass (\S\ref{sec:commit-reveal}) does nothing here,
because none of the public writes are concurrent with each
other.

\begin{center}
\begin{tikzpicture}[
  every node/.style={font=\small},
  ev/.style={draw, rounded corners, inner sep=2pt, minimum width=2.1cm,
             minimum height=4mm, fill=blue!6},
  cnode/.style={draw, inner sep=2pt, minimum width=2.1cm, minimum height=4mm,
             fill=red!10},
  rv/.style={draw, regular polygon, regular polygon sides=6,
             inner sep=0pt, minimum width=2.3cm, minimum height=5mm,
             fill=green!10},
  arr/.style={-{Stealth[length=2mm]}, thick},
]
\node[ev]    (j0) at (0,0)   {join Host};
\node[ev]    (j1) at (0,-1)  {join Guest};
\node[cnode] (c2) at (0,-2)  {commit Host.car};
\node[ev]    (y3) at (0,-3)  {yield Guest.d};
\node[ev]    (y4) at (0,-4)  {yield Host.goat};
\node[ev]    (y5) at (0,-5)  {yield Guest.switch};
\node[rv]    (r6) at (0,-6)  {reveal Host.car};

\draw[arr] (j0) -- (j1);
\draw[arr] (j1) -- (c2);
\draw[arr] (c2) -- (y3);
\draw[arr] (y3) -- (y4);
\draw[arr] (y4) -- (y5);
\draw[arr] (y5) -- (r6);
\draw[arr, bend left=55] (c2.east) to (r6.east);
\end{tikzpicture}
\end{center}

Two facts to read off this picture.  First, the long-distance
edge from \TT{commit Host.car} to \TT{reveal Host.car} captures
the commit-reveal pairing structurally: the reveal is enabled
only after every preceding move, including its own committed
predecessor.  In any configuration where the reveal is on the
frontier, every prior decision has been made.  Second, the data
dependency from \TT{yield Guest.d} to \TT{yield Host.goat} ---
because Host's guard reads \TT{Guest.d} --- is collapsed into
the phase-order edge, so it does not appear separately; the
guard read is redundant given the existing chain.

This is the typical shape of turn-taking games: a sequential
chain plus the structural pair for hidden information.
Concurrency-driven rewriting becomes interesting when several
roles share a single yield statement, as in the Vickrey example
treated in \S\ref{sec:commit-reveal}.

\subsection{Visibility-aware perfect recall}
\label{sec:eventgraph:visibility-recall}

For game-theoretic analysis, two players are in the same
information set at a node $B$ iff they have observed the same
sequence of public events leading to $B$.  In the EventGraph this
becomes a structural notion: the set of fields observable by role
$R$ at node $B$ is

\begin{equation*}
  \mathrm{obs}_R(B) \;=\; \bigl\{\, f \;\big|\; f \text{ is written by some predecessor } A \text{ of } B,
            \text{ and either } \mathrm{owner}(A) = R \text{ or } \mathrm{vis}(A,f) \in \{\PUBLIC, \REVEAL\} \bigr\}.
\end{equation*}

This is the projection used by the Gambit backend to compute
information sets and, in a closely related form, by the MAID and
Scribble backends.  An IR helper \TT{observableFieldsAt(B)}
computes it; the formula above is the only place where
observability is defined.

\subsection{Implementation notes}

The EventGraph class wraps a generic \TT{Dag<NodeId>} with the
per-node payload map and a precomputed reachability relation.
Reachability is computed once when the graph is built and reused
for all queries; building it eagerly is the right tradeoff in
this setting because the graph is small (forty-some nodes for
the largest example) and the queries are many
(\TT{reaches}, \TT{ancestorsOf}, \TT{canExecuteConcurrently}).

The constructor takes a node set, a per-node predecessor map, and
a per-node payload map, and returns either an EventGraph or
\TT{null}; the latter happens when an invariant is violated.
Three internal validators run inside the constructor:

\begin{itemize}\itemsep1pt
\item acyclicity, via topological sort;
\item commit-reveal ordering, by checking that every \REVEAL has
  a matching \COMMIT among its ancestors;
\item read-visibility, by checking that every field appearing in a
  guard at node $B$ is either written \PUBLIC before $B$, or has
  been revealed before $B$, or is being written by $B$ itself.
\end{itemize}

The corresponding error in the frontend is a static error, with a
source span pointing back to the offending \TT{where} clause or
\TT{commit/reveal} pair.  This is the same set of checks a
human reviewer would perform when reading a hand-written
commit-reveal contract --- the IR just performs them
systematically.

A note on the read-visibility check: in addition to fields written
\PUBLIC or revealed by a predecessor, the validator also accepts a
guard read of a field that the actor itself committed earlier on
the same DAG path.  An owner ``sees their own writes'' is therefore
a slightly stronger statement than the bullet above admits: it
includes their own \emph{committed} fields as well as the
public/revealed ones.  In practice the typechecker's pre-IR
desugaring strips fields written by the current node from the
guard's scope, so this branch matters mostly for guards that
reference an actor's earlier private writes.
